\chapter{When order wins}\label{ch:oc-bias}

This chapter collects three places where a random model is subtly wrong: Chebyshev's bias, Maier's irregularities, and the bias of Lemke Oliver and Soundararajan for consecutive primes.

\section{Chebyshev's bias}

Chebyshev observed that, although the prime-number theorem for arithmetic progressions gives $\pi(x;4,3)\sim\pi(x;4,1)$, the difference $\pi(x;4,3)-\pi(x;4,1)$ is positive for most small $x$. Rubinstein and Sarnak explained the logarithmic density of the lead under the linear independence of the zeros.

\section{Maier's irregularities}

Maier's matrix method produces short intervals that contain more primes than the Cram\'er model allows, and others that contain fewer. The irregularities are deterministic and repeat at every scale $(\log x)^\lambda$.

\section{Consecutive primes}

Lemke Oliver and Soundararajan found that consecutive primes avoid repeating a residue class more than a random model predicts. The bias is visible in the singular series of the Hardy--Littlewood conjecture, and it is a theorem under hypotheses on the zeros.
